\begin{proposition}[6.11.6]
\label{IV.2.s6.11.6-fr}
Sean $X$ un preesquema localmente noetheriano, $\mathcal F$ un
$\mathcal O_X$-Módulo coherente y $n$ un entero $>0$. Supongamos que, para todo
subpreesquema cerrado íntegro $Y$ de $X$, existe una parte abierta no
vacía $W$ de $Y$ tal que el subpreesquema de $Y$ inducido sobre el abierto $W$
verifica $(S_n)$. Entonces el conjunto $U_{S_n}(\mathcal F)$ es abierto en
$X$.
\end{proposition}

\oldpage[IV-2]{161}

\begin{proof}
Como la cuestión es local en $X$, es posible limitarse al caso en que $X$ es
noetheriano. Razonaremos por inducción sobre $n$: para $n=1$,
el conjunto $U_{S_1}(\mathcal F)$ es abierto en $X$, pues el conjunto de los
$x\in X$ en los que $\mathcal F$ no verifica $(S_1)$ es el conjunto de los $x$ tales
que $\mathcal F_x$ admite ideales primos asociados inmersos (5.7.5);
si $(Z_j)$ es la familia finita de los ciclos primos asociados a $\mathcal F$
que son inmersos, se tiene, por tanto,
$U_{S_1}(\mathcal F)=X\setmin\bigcup_jZ_j$, lo que prueba nuestra afirmación, ya que los
$Z_j$ son cerrados. Supondremos, pues, en adelante que $n>1$. En segundo
lugar, es posible limitarse al caso en que
$\operatorname{Supp}(\mathcal F)=X$, pues existe un subpreesquema cerrado
$T$ de $X$ cuyo espacio subyacente es $\operatorname{Supp}(\mathcal F)$,
y un $\mathcal O_T$-Módulo coherente $\mathcal G$ tales que
$\mathcal F=j_*(\mathcal G)$, siendo $j:T\to X$ la inyección canónica
(\textbf{Err}$_{\rm III}$, 30). Como es equivalente probar que $\mathcal F$
o $\mathcal G$ verifica $(S_n)$ en un punto
$x\in\operatorname{Supp}(\mathcal F)$, basta considerar el
caso en que $T=X$. Observemos, por último, que, por definición (5.7.2),
$U_{S_n}(\mathcal F)$ es estable por generización. Utilizaremos el
lema siguiente:

\begin{lemma}[6.11.6.1]
\label{IV.2.s6.11.6.1-fr}
Sean $X$ un espacio noetheriano en el que toda parte cerrada irreducible
admite un punto genérico, y $E$ una parte de $X$. Para que $E$ sea abierta
en $X$, es necesario y suficiente que $E$ sea estable por generización y que,
para toda parte abierta $V$ de $X$ y toda parte irreducible $Y$ cerrada
en $V$, tales que $V\setmin Y\subset E$ y el punto genérico de $Y$
pertenece a $E$, $E\cap Y$ contenga una parte abierta no vacía de $Y$.
\end{lemma}

\begin{proof}
Se observará que este criterio implica el de (\textbf{0}$_{\rm III}$, 9.2.6)
cuando toda parte cerrada irreducible de $X$ admite un punto genérico.
Evidentemente, solo hay que demostrar la suficiencia de las condiciones enunciadas.

Consideremos el interior $U$ de $E$; el conjunto cerrado $X\setmin U$ es la unión de
sus componentes irreducibles, que son finitas en número y cerradas en $X$.
Si se tuviese $E\ne U$, la hipótesis de que $E$ es estable por generización
implicaría que el punto genérico $z$ de una de las componentes irreducibles
de $X\setmin U$ pertenecería a $E$. Ahora bien, $z$ pertenece a una sola de las
componentes irreducibles de $X\setmin U$; si $T$ es la unión de las otras
componentes irreducibles de $X\setmin U$, $V=X\setmin T$ es abierto en $X$ y es la unión de
$U$ y del conjunto $Y=\overline{\{z\}}\cap V$, cerrado en $V$ e
irreducible. Por hipótesis, $E\cap Y$ contiene una parte $W$ abierta \emph{en}
$Y$; se deduce que $U\cup W$ es abierto \emph{en} $V$ y, por tanto, en $X$, lo que
es absurdo, ya que se ha supuesto que $U$ es el interior de $E$.
\end{proof}

En virtud de este lema, se puede suponer que existe en $X$ un
subpreesquema íntegro $Y$ con punto genérico $y$ tal que $\mathcal F$
verifica $(S_n)$ en el punto $y$ y en todos los puntos de $X\setmin Y$, y todo se
reduce a comprobar que existe en $X$ una vecindad abierta de $y$ tal que
$\mathcal F$ verifica $(S_n)$ en esa vecindad. Distingamos dos casos:

$1^\circ$ $y$ es un punto maximal de $X$; como existe una vecindad abierta
de $y$ que no encuentra ninguna componente irreducible de $X$ distinta de
$\overline{\{y\}}$, se puede suponer que $X$ es irreducible y, por tanto, tiene el mismo
espacio subyacente que $Y$, de modo que $Y$ está definido por el Nilradical de
$\mathcal O_X$, que es \emph{nilpotente}. Por otra parte, reemplazando $X$
por una vecindad abierta de $y$, se puede suponer que $\mathcal F$ es \emph{normalmente
plano} a lo largo de $Y$ (6.9.1); se sigue entonces de (6.10.4) que
$U_{S_n}(\mathcal F)=U_{S_n}(\mathcal O_Y)$, y como este último es, por
hipótesis, una vecindad de $y$ en $X$, ello prueba la proposición en este
caso.

\oldpage[IV-2]{162}

$2^\circ$ $y$ no es un punto maximal de $X$, es decir (puesto que
$\operatorname{Supp}(\mathcal F)=X$), $\dim(\mathcal F_y)\geq1$, y, por tanto, también
$\operatorname{prof}(\mathcal F_y)\geq1$, ya que $\mathcal F$ verifica por
hipótesis $(S_n)$ (y \emph{a fortiori} $(S_1)$) en el punto $y$. Reemplazando, si es
necesario, $X$ por una vecindad abierta de $y$, se puede, por tanto, suponer
que existe una sección $f$ de $\mathcal O_X$ sobre $X$, $\mathcal F$-regular y
tal que $f_y\in\mathfrak m_y$, o, equivalentemente, $f(y)=0$ (\textbf 0, 15.2.4); se tiene,
por consiguiente, $f(x)=0$ para todo $x\in Y$ (\textbf{0}$_{\rm I}$, 5.5.2). Se sabe que
$\mathcal F/f\mathcal F$ verifica $(S_{n-1})$ en el punto $y$ (5.7.6).
Aplicando la hipótesis de inducción y reemplazando, si es necesario, $X$ por una
vecindad abierta de $y$, se puede, por tanto, suponer que
$\mathcal F/f\mathcal F$ verifica $(S_{n-1})$ en todo punto de $X$. Pero, para
todo $x\in Y$, la relación $f(x)=0$ implica que se tiene
$\operatorname{prof}(\mathcal F_x/f_x\mathcal F_x)=
\operatorname{prof}(\mathcal F_x)\setmin 1$ y
$\dim(\mathcal F_x/f_x\mathcal F_x)=\dim(\mathcal F_x)\setmin 1$
(\textbf 0, 16.3.4 y 16.4.6); por tanto, la relación
\[
\operatorname{prof}(\mathcal F_x/f_x\mathcal F_x)\geq
\inf(n\setmin 1,\dim(\mathcal F_x/f_x\mathcal F_x)
\]
equivale a
\[
\operatorname{prof}(\mathcal F_x)\geq\inf(n,\dim(\mathcal F_x)).
\]
Como se ha supuesto que $\mathcal F$ verifica $(S_n)$ en todo punto de $X\setmin Y$,
esto concluye la demostración.
\end{proof}
